\begin{abstract}
Do Large Language Models (LLMs) and biological brains converge onto similar computational mechanisms during language processing? Applying Marr’s tri-level framework reveals substantial mismatches between humans and LLMs at both the computational level (broad communicative intent versus next-token prediction) and the implementational level (biological neural circuits versus artificial Transformer/SSM architectures). Whether these systems nevertheless converge at the algorithmic level -- the representations and step-by-step processing dynamics -- remains an open question. To test for algorithmic alignment, cognitive scientists have increasingly developed "human-like" models constrained by human cognitive limitations (e.g., working memory bottlenecks) or developmentally plausible training exposure (e.g., the BabyLM challenge). Here, we evaluate algorithmic convergence by establishing a multi-modal human benchmark, evaluating how accurately model-derived metrics (surprisal and internal contextual embeddings) predict human behavioral responses (reading times and eye-tracking fixations) and brain activity (fMRI, MEG, EEG event-related potentials, and intracranial ECoG recordings). Counter to the hypothesis that human-like inductive biases guarantee greater cognitive alignment, our results demonstrate that current human-like models do not outperform standard, unconstrained LLMs in predicting human neural and behavioral metrics at the algorithmic level. Instead, the strong alignment observed in standard LLMs despite implementational differences lends empirical support to the Platonic Representation Hypothesis, suggesting that task-scale optimization naturally drives disparate information-processing systems toward shared representations. We make available our multi-modal evaluation framework as the LLM-Human Convergence Benchmark, providing a systematic testing ground to guide the future design of genuinely human-aligned cognitive architectures.

\end{abstract}